\documentclass[10pt]{beamer}

%\documentclass[10pt, handout]{beamer}
\setbeameroption{show notes}

%\documentclass[10pt, a4paper]{article}
%\usepackage{beamerarticle}

\input{preambule.tex}


\usepackage{media9}

\graphicspath{{Images/}{Images/\jobname/}}

\date{\today}


\renewcommand{\LectionNumber}{14}
\renewcommand{\LectionTheme}{Основные схемы представления объемных тел}
\title{Лекция \lecdni \\ \LectionTheme}
\subtitle{Компьютерная графика}



%\usepackage{standalone}

\setbeamersize
{
	text margin left=0.5cm,
	text margin right=0.5cm
}

\usepackage{comment}


%	\transduration{2}
%   \transfade

\begin{document}
 		 
\makeatletter
\defbeamertemplate*{title page}{my theme}
{
	
	\hfill
	
	\begin{beamercolorbox}[wd=.9\paperwidth,center,]{title}%
		
	\end{beamercolorbox}%	
	
	\vbox to 1em {}
	
	\begin{beamercolorbox}[wd=.9\paperwidth,center,]{title}%
		\usebeamerfont{subtitle}%
		\hfill
		
		\insertsubtitle
		
		\usebeamerfont{title}%
		\inserttitle{} \\[0.5em]
		
		
		
	\end{beamercolorbox}%	
	\hfill\hfill
	
	\begin{beamercolorbox}[wd=.9\paperwidth,center,]{}%
		\usebeamerfont{author}%
		\hfill \\[0.5em]
		\insertauthor{} \\[0.5em]
		\regals
		    
		\vbox to 1em{}
		\usebeamerfont{institute}%
		\insertinstitute {}
		
		\vbox to 1em{}			
		
	\end{beamercolorbox}%	
	\hfill\hfill
	
	
	
	{ \usebeamerfont{institute}\hfill Обновлено: \insertshortdate{} }
	
	\vbox to 5em{}
	
	
}
\defbeamertemplate*{footline}{my theme}{
	\leavevmode%
	\hbox{%
		\begin{beamercolorbox}[wd=.25\paperwidth,ht=3.25ex,dp=0ex,center,sep=1pt]{author in head/foot}%
			\usebeamerfont{author in head/foot}%
			\insertauthor 
			\beamer@ifempty{\insertshortinstitute}{}
		\end{beamercolorbox}%
		\begin{beamercolorbox}[wd=.65\paperwidth,ht=3.25ex,dp=0ex,center,sep=1pt]{title in head/foot}%
			\usebeamerfont{title in head/foot}\insertshortinstitute
		\end{beamercolorbox}%
		\begin{beamercolorbox}[wd=.1\paperwidth,ht=3.25ex,dp=0ex,center, sep=0.5pt]{date in head/foot}%
			\usebeamerfont{date in head/foot}
			\footnotesize \expandafter\toDni{\insertframenumber} / \expandafter\toDni{\inserttotalframenumber}
	\end{beamercolorbox}}%
}

\makeatletter
\beamer@theme@subsectionfalse
\makeatother

 % Head
%\defbeamertemplate*{headline}{miniframes theme}
%{%
%	\begin{beamercolorbox}[colsep=1.5pt]{upper separation line head}
%	\end{beamercolorbox}
%	\begin{beamercolorbox}{section in head/foot}
%		\vskip2pt\insertnavigation{\paperwidth}\vskip2pt
%	\end{beamercolorbox}%
%	\ifbeamer@theme@subsection%
%	\begin{beamercolorbox}[colsep=1.5pt]{middle separation line head}
%	\end{beamercolorbox}
%	\begin{beamercolorbox}[ht=2.5ex,dp=1.125ex,%
%		leftskip=.3cm,rightskip=.3cm plus1fil]{subsection in head/foot}
%		\usebeamerfont{subsection in head/foot}\insertsubsectionhead
%	\end{beamercolorbox}%
%	\fi%
%	\begin{beamercolorbox}[colsep=1.5pt]{lower separation line head}
%	\end{beamercolorbox}
%}

{
	
}


\makeatother






%float page top aligment
\makeatletter
\setlength{\@fptop}{0pt}
\setlength{\@fpbot}{0pt plus 1fil}
\makeatother

\newcommand \abs[1] {\left| #1 \right|}

\everymath{\displaystyle}

    
\begin{comment}
\end{comment}


    
    
    \QRFRAME	
	

	\frame{\maketitle}

	
	\begin{frame}\frametitle{План лекции}
		\tableofcontents
	\end{frame}
	




%\section{Рациональные кривые}
%\frame {\sectionpage}


\begin{frame}\frametitle{Зачем моделировать твердые тела?}
			
	\begin{columns}
		\begin{column}{0.33\textwidth}
			\centering
			\IG{tflex.jpg}
			
			CAD
		\end{column}
		\begin{column}{0.33\textwidth}
			\centering
			\IG{chpu.jpg}
			
			Станки с ЧПУ
		\end{column}
		
		\begin{column}{0.33\textwidth}
			\centering
			\IG{r6.jpg}
			
			Роботы
		\end{column}
	\end{columns}	
	
		\begin{columns}
		\begin{column}{0.5\textwidth}
			\centering
			\IG[0.7]{preview-blog-42-01.jpg}
			
			3D - печать
		\end{column}
		\begin{column}{0.5\textwidth}
			\centering
			\IG[0.5]{Screenshot 2025-10-23 162528.png}
			
			CGI
		\end{column}
		
		
	\end{columns}	
	
\end{frame} 




\section[B-Rep]{Гранчиное представление }

\begin{frame}\COVER{\frametitle{Граничное представление (boundary representation, b-rep)}}

	
	\begin{block}{Определение}
		Твёрдое тело --- это замкнутое пространство, ограниченное множеством поверхностей. При этом важно, что проводится различие между внутренностью объекта, “внешним” пространством и поверхностями, ограничивающими тело. 
	\end{block}
	
	\begin{center}
		\IG[0.7]{Screenshot 2025-10-30 155040.png}
	\end{center}
	
\end{frame}


\begin{frame}\COVER{\frametitle{}}
	
	\TC{0.5}
	{
		\IG{Screenshot 2025-10-23 164125.png}
	}{
		Многогранник - это твёрдое тело, ограниченное множеством полигональных граней, чьи ребра принадлежат четному числу граней. 
		
		Простой многогранник -   может быть деформировани в сферу. 
		
		b-rep простого многранника удовлетворяет формуле Эйлера:
		
		$V - E + F = 2$.
		
	}
	
	
\end{frame}

\begin{frame}\COVER{\frametitle{}}
	
	\IG{Screenshot 2025-10-30 145748.png}
	
\end{frame}

\begin{frame}\COVER{\frametitle{}}
	
	\TCT{0.5}
	{
		\textbf{Преимущества}
		
		\begin{itemize}
			\item Высокая точность и детализация.
			
			\item Прямая пригодность для рендеринга.
			
			\item Стандарт в CAD-системах (SolidWorks, Fusion 360).
			
		\end{itemize}
	}
	{
		\textbf{Недостатки}
		
		\begin{itemize}
			
			\item Сложность булевых операций.
			
			\item Возможность некорректных моделей ( <<дыры>>, неориентированные грани).
			
			\item Большой объем памяти для сложных моделей.
			
		\end{itemize}
	}
	
\end{frame}

\section[CSG]{Конструктивная блочная геометрия (CSG)}

\begin{frame}\COVER{\frametitle{Конструктивная блочная геометрия (CSG)}}
	
	\TC{0.4}
	{
		\IG{Csg_tree.png}
		
		

	}
	{
		
		\centering 
		
		\TCT{0.5}
		{
			\centering
			
			\IG{Boolean_union.png} 
			
			$C \cup S$  
		}
		{
			\centering
			
			\IG{Boolean_intersect.png}
			
			$C \cap S$  
		}
		
		\IG[0.5]{Boolean_difference.png}
		
		$ C \setminus  S$  (разность)
		
	}
	
	
	\begingroup 
	\renewcommand\thefootnote{}
	\footnotetext{\HREF{https://www.youtube.com/watch?v=Ka6vjgrDeiU}{Honda CBX 1000 Motorcycle Engine}}
	\endgroup
	
	
\end{frame}



\begin{frame}\COVER{\frametitle{Преимущества и Недостатки}}
	
\centering
	
		\TCT{0.5}
		{
			\textbf{Преимущества}
			
			\begin{itemize}
				\item Компактность хранения (только дерево). 
				\item Параметричность и легкость редактирования.
				\item Гарантированная корректность (для простых примитивов).
			\end{itemize}
		}
		{
			\textbf{Недостатки}
			
			\begin{itemize}
				\item Неоднозначность представления.
				\item Для отображения нужно конвертировать в B-Rep или полигоны.
				\item Ограниченный набор форм "из коробки".
			\end{itemize}
		}
	
	
	\IG[0.45]{csg1.jpg}
	
\end{frame}

\section[Воксели]{Воксельное представление}


\begin{frame}\COVER{\frametitle{Воксельное представление}}
	
	\TC{0.5}
	{
		\IG{Voksel-kirpich-dlya-trehmernoj-grafiki-768x449.jpg}
	}
	{
	}
	
	
	\begingroup
	\renewcommand\thefootnote{}
	\footnotetext{\HREF{https://www.youtube.com/watch?v=8ptH79R53c0}{New Voxel Engine Reveal - Crystal Islands Experiment}}
	\footnotetext{\HREF{https://iolite-engine.com/}{IOLITE | Voxel Game Engine}}
	\endgroup
	
	
\end{frame}

\section[Octrees]{Восьмиричные деревья}


\begin{frame}\COVER{\frametitle{Восьмиричные деревья}}
	
	\IG{7438d4db9a93c3bf5caf8148b6562337.png}
	
\end{frame}

\begin{frame}\COVER{\frametitle{Поиск коллизий между объектами}}

	\centering\IG[0.7]{2b948a088e81982d6eb572cc24d866c2.png}
	
	
\end{frame}


\end{document}


